\worksheettitle
  {Ломаная и многоугольник}
  {\modulemark{M05} Учебный модуль \textbullet\ полный маршрут}

\studentnameblock

\begin{notebox}[Чему вы научитесь]
  Различать ломаную и многоугольник, определять замкнутость и правильно
  называть вершины, звенья и стороны.
\end{notebox}

\begin{checkpointbox}[Вход: соберите фигуру]
  Отрезки $AB$, $BC$ и $CA$ соединены концами.
  \begin{enumerate}[label=\textbf{\arabic*.}]
    \item Полученная граница замкнута? \answerblank[20mm]
    \item Сколько у неё вершин? \answerblank[15mm]
    \item Как называется фигура? \answerblank[30mm]
  \end{enumerate}
\end{checkpointbox}

\begin{notebox}[Как устроен маршрут]
  Сначала разберите устройство ломаной. Затем сравните замкнутые и
  незамкнутые ломаные, выделите границу многоугольника, потренируйтесь
  называть его элементы, исправьте ошибку и выполните контроль.
\end{notebox}

\section{Разбираем ломаную}

\openbrokenfigure

\begin{fillbox}[Что видим на рисунке]
  Ломаная $ABCDE$ состоит из последовательно соединённых отрезков.
  Точки $A,B,C,D,E$ --- её \answerblank[25mm], а отрезки
  $AB,BC,CD,DE$ --- её \answerblank[25mm].

  Здесь \answerblank[15mm] вершин и \answerblank[15mm] звена.
\end{fillbox}

\begin{notebox}[Правило]
  Ломаная состоит из отрезков, соединённых последовательно: конец одного
  звена является началом следующего. Если последний конец совпадает с первым,
  ломаная замкнута; иначе она незамкнута.
\end{notebox}

\section{Практика с опорой}

\studentrecognizefigure

\begin{practicebox}[1. Определите тип]
  \begin{enumerate}[label=\alph*)]
    \item Незамкнутые ломаные: \answerblank[30mm].
    \item Замкнутые ломаные: \answerblank[30mm].
    \item На рисунке 1 число вершин \answerblank[12mm], число звеньев
      \answerblank[12mm].
  \end{enumerate}
\end{practicebox}

\section{От замкнутой ломаной к многоугольнику}

\polygonpartsfigure

\begin{fillbox}[Разберите пример]
  Граница шестиугольника $ABCDEF$ --- замкнутая ломаная.
  Её звенья называют \answerblank[25mm] многоугольника, а их концы ---
  \answerblank[25mm]. У шестиугольника $6$ сторон и $6$ вершин.
\end{fillbox}

\polygonnamesfigure

\begin{notebox}[Название по числу сторон]
  Три стороны --- треугольник, четыре --- четырёхугольник, пять ---
  пятиугольник, шесть --- шестиугольник. У многоугольника число сторон
  равно числу вершин.
\end{notebox}

\begin{practicebox}[2. Назовите элементы]
  Для шестиугольника $ABCDEF$ на рисунке запишите:
  \begin{enumerate}[label=\alph*)]
    \item вершины: \answerblank[95mm];
    \item стороны: \answerblank[95mm];
    \item число сторон и вершин: \answerblank[50mm].
  \end{enumerate}
\end{practicebox}

\section{Самостоятельная практика}

\begin{practicebox}[3. Проверьте описание]
  Замкнутая ломаная $KLMNK$ образует многоугольник.
  \begin{enumerate}[label=\alph*)]
    \item Назовите многоугольник по числу сторон. \answerblank[32mm]
    \item Запишите все его стороны. \answerblank[85mm]
    \item Объясните, почему $KM$ не является стороной этой границы.
      \writinglines{2}
  \end{enumerate}
\end{practicebox}

\begin{practicebox}[4. Постройте]
  Отметьте четыре точки $P,Q,R,S$ и соедините их так, чтобы получилась
  незамкнутая ломаная из трёх звеньев. Затем добавьте одно звено, чтобы
  ломаная стала границей четырёхугольника.
  \writinglines{4}
\end{practicebox}

\section{Разбираем ошибку}

\begin{warningbox}[«Любая замкнутая линия --- многоугольник»]
  Ученик назвал окружность многоугольником, потому что она замкнута.

  \textbf{Причина ошибки:} \answerblank[100mm]

  \textbf{Исправление:} \answerblank[100mm]
\end{warningbox}

\section{Контрольная точка}

\openclosedfigure

\begin{checkpointbox}[5. Проверьте результат]
  \begin{enumerate}[label=\textbf{\arabic*.}]
    \item Какая ломаная замкнута? \answerblank[30mm]
    \item Назовите вершины и звенья незамкнутой ломаной. \writinglines{2}
    \item Сколько сторон и вершин у многоугольника справа?
      \answerblank[35mm]
    \item Назовите его, перечисляя вершины подряд по границе.
      \answerblank[45mm]
  \end{enumerate}
\end{checkpointbox}

\begin{reflectionbox}[Проверяю себя]
  \begin{itemize}[label=\choicebox]
    \item различаю вершины и звенья ломаной;
    \item определяю, замкнута ли ломаная;
    \item выделяю границу многоугольника;
    \item называю стороны и вершины по порядку.
  \end{itemize}
\end{reflectionbox}

\begin{resultbox}[Дальнейший маршрут]
  Если путаете вершины, звенья, стороны или замкнутость, выполните M05-R.
  Для тренировки используйте M05-H. При устойчивом результате переходите
  к M06. M05\(\star\) выполняется дополнительно.
\end{resultbox}
